%=============================================================================
% AGENT 5: PULSED OPERATION FOR ATMOSPHERIC PSI-BUBBLE PROPULSION
% The timing solution: use pulsed magnetic fields to simultaneously evacuate
% gas AND cross threshold before gas can return.
%
% DFD parameters:
%   n = e^psi,   a = (c^2/2) grad(psi),   kappa = 3/(8 alpha) ~ 51.4,
%   eta_c = alpha/4 ~ 1.82e-3
%=============================================================================

\documentclass[amsmath,amssymb,aps,prd,superscriptaddress,nofootinbib]{revtex4-2}
\usepackage{amsmath,amssymb,bm,graphicx,physics,siunitx,booktabs}
\usepackage{hyperref}

\newcommand{\psifield}{\psi}
\newcommand{\avec}{\mathbf{a}}
\newcommand{\ka}{\kappa}
\newcommand{\etac}{\eta_c}

\begin{document}

\title{Pulsed Magnetic Field Operation for Atmospheric $\psi$-Bubble Propulsion:\\
Gas Evacuation, Threshold Crossing, and Impulse Analysis}

\author{DFD Research Programme -- Agent 5 (Pulsed Operation)}
\date{28 March 2026}

\begin{abstract}
The central obstacle to atmospheric $\psi$-bubble propulsion is that
sea-level air density ($\rho = 1.225$~kg/m$^3$) keeps the DFD
threshold ratio $\eta = U_\text{EM}/(\rho c^2)$ far below the critical
value $\etac = \alpha/4 \approx 1.82 \times 10^{-3}$. We show that
\emph{pulsed} magnetic fields offer a pathway: a high-$B$ pulse
($\sim$50~T) generates magnetic pressure exceeding $10^4$~atmospheres,
evacuating the near-field region in $\sim$5~$\mu$s---faster than air can
return. During the brief evacuated window, $\eta$ crosses threshold and
the $\kappa$-channel activates, accelerating residual gas to $\sim
0.01\,c$. We compute the full time-dependent density profile, threshold
crossing dynamics, impulse per pulse, average thrust at various
repetition rates, and power requirements. We conclude that while
full-bore 50~T pulsed operation at 10~kHz is power-prohibitive ($\sim$10~GW),
an optimised scheme modulating between a persistent 18~T base field and
brief 22~T pulses at 100~Hz requires $\sim$64~MW---achievable with a
compact nuclear reactor---and produces meaningful thrust for atmospheric
flight.
\end{abstract}

\maketitle

%=============================================================================
\section{The Atmospheric Threshold Problem}\label{sec:problem}
%=============================================================================

\subsection{Why Atmosphere Matters}

In vacuum or near-vacuum, the DFD threshold condition
\begin{equation}\label{eq:threshold}
\eta \equiv \frac{U_\text{EM}}{\rho_m c^2} \geq \etac = \frac{\alpha}{4}
\approx 1.82 \times 10^{-3}
\end{equation}
is easily satisfied: with $\rho_m \to 0$, any finite EM field
gives $\eta \to \infty$. The vacuum $\psi$-bubble analysis (Agent~4)
showed that even modest ferrite-level fields ($B_\text{surface} \sim
3.5 \times 10^{-3}$~T) create above-threshold regions extending to
$\sim$8~m at sea level and hundreds of metres in moderate vacuum.

At sea level, however, $\rho_\text{air} = 1.225$~kg/m$^3$ and the
threshold EM energy density required is:
\begin{equation}
U_\text{EM}^\text{thresh} = \etac \, \rho_\text{air} \, c^2
= 1.82 \times 10^{-3} \times 1.225 \times 9 \times 10^{16}
= 2.01 \times 10^{14} \; \si{J/m^3}.
\end{equation}

The corresponding magnetic field:
\begin{equation}
B_\text{thresh} = \sqrt{2\mu_0 \, U_\text{EM}^\text{thresh}}
= \sqrt{2 \times 4\pi \times 10^{-7} \times 2.01 \times 10^{14}}
= \sqrt{5.05 \times 10^{8}}
\approx 22{,}500 \; \si{T}.
\end{equation}

This is roughly 500 times beyond the capability of any terrestrial
magnet---steady-state atmospheric threshold crossing is not feasible
through brute-force field strength alone.

\subsection{The Pulsed Operation Concept}

The key insight is that the threshold depends on the \emph{local}
density $\rho_m$, not the ambient density $\rho_\text{air}$. If we
can transiently reduce $\rho_m$ in the near-field region, we can cross
threshold with achievable field strengths.

A strong pulsed magnetic field provides exactly this mechanism:
\begin{enumerate}
\item \textbf{Magnetic pressure evacuates gas.} A 50~T field exerts
  $B^2/(2\mu_0) = 10^9$~Pa $\approx 10{,}000$~atm, accelerating air
  outward at $\sim 10^{10}$~m/s$^2$.
\item \textbf{Gas clears in microseconds.} The evacuation time for a
  0.1~m region is $\sim$5~$\mu$s.
\item \textbf{Threshold is crossed in the evacuated region.} With
  $\rho_m \ll \rho_\text{air}$, the threshold field drops dramatically.
\item \textbf{$\kappa$-coupling activates.} Residual gas in the
  above-threshold zone is accelerated to high velocity.
\item \textbf{Field relaxes, gas returns.} Over $\sim$100~$\mu$s, air
  refills the evacuated region.
\item \textbf{Repeat.} At a suitable repetition rate, sustained
  average thrust is generated.
\end{enumerate}

%=============================================================================
\section{Gas Dynamics During a 50~T Pulse}\label{sec:gas-dynamics}
%=============================================================================

\subsection{Magnetic Pressure and Initial Acceleration}

The magnetic pressure from a 50~T field:
\begin{equation}\label{eq:Pmag}
P_\text{mag} = \frac{B^2}{2\mu_0}
= \frac{(50)^2}{2 \times 4\pi \times 10^{-7}}
= \frac{2500}{2.513 \times 10^{-6}}
= 9.95 \times 10^{8} \; \si{Pa}
\approx 10^9 \; \si{Pa}.
\end{equation}

This exceeds atmospheric pressure by a factor of $\sim 10^4$.

The initial acceleration of air at the magnetic field gradient
(gradient length $L \sim 0.1$~m):
\begin{equation}\label{eq:a-gas}
a_\text{gas} = \frac{P_\text{mag}}{\rho_\text{air} \cdot L}
= \frac{10^9}{1.225 \times 0.1}
= 8.16 \times 10^{9} \; \si{m/s^2}.
\end{equation}

For comparison, this is $\sim 8.3 \times 10^{8}$~g---an extraordinary
acceleration that will move gas at nearly explosive speed.

\subsection{Time-Dependent Density Profile}

We model the gas evacuation as one-dimensional motion of a gas element
at initial position $x_0$ from the magnet surface, subject to
acceleration $a_\text{gas}$. The position of the gas front is:
\begin{equation}
x(t) = x_0 + \tfrac{1}{2} a_\text{gas} \, t^2.
\end{equation}

The time for the gas front to travel a distance $L$ from its initial
position:
\begin{equation}\label{eq:t-clear}
t_\text{clear} = \sqrt{\frac{2L}{a_\text{gas}}}
= \sqrt{\frac{2 \times 0.1}{8.16 \times 10^9}}
= \sqrt{2.45 \times 10^{-11}}
= 4.95 \times 10^{-6} \; \si{s}
\approx 5 \; \mu\text{s}.
\end{equation}

The density as a function of time and position follows from mass
conservation. For a gas element initially at position $x_0$ with
density $\rho_0$:
\begin{equation}
\rho(x_0, t) = \rho_0 \cdot \frac{dx_0}{dx}
= \frac{\rho_0}{1 + a_\text{gas}\,t^2/(2x_0)}.
\end{equation}

Near the magnet surface ($x_0 \to 0$), the density drops as:
\begin{equation}\label{eq:rho-t}
\rho(0, t) \approx \rho_0 \cdot \frac{L}{L + \tfrac{1}{2}a_\text{gas}\,t^2}
\end{equation}
where we have smoothed the singularity by treating the initial gas
column of length $L$ as a slab.

\begin{center}
\begin{tabular}{cccc}
\toprule
$t$ ($\mu$s) & $\frac{1}{2}a\,t^2$ (m) & $\rho/\rho_0$ & $\rho$ (kg/m$^3$) \\
\midrule
0 & 0 & 1 & 1.225 \\
1 & $4.1 \times 10^{-3}$ & 0.96 & 1.18 \\
2 & $1.6 \times 10^{-2}$ & 0.86 & 1.05 \\
3 & $3.7 \times 10^{-2}$ & 0.73 & 0.89 \\
4 & $6.5 \times 10^{-2}$ & 0.61 & 0.74 \\
5 & $1.0 \times 10^{-1}$ & 0.50 & 0.61 \\
7 & $2.0 \times 10^{-1}$ & 0.33 & 0.41 \\
10 & $4.1 \times 10^{-1}$ & 0.20 & 0.24 \\
15 & $9.2 \times 10^{-1}$ & 0.098 & 0.12 \\
20 & $1.6$ & 0.059 & 0.072 \\
30 & $3.7$ & 0.026 & 0.032 \\
50 & $10.2$ & 0.0097 & 0.012 \\
70 & $20$ & 0.0050 & 0.0061 \\
100 & $41$ & 0.0024 & 0.0030 \\
\bottomrule
\end{tabular}
\end{center}

\subsection{Refinement: Compressible Flow}

The above simple model neglects the finite sound speed ($c_s \approx
343$~m/s in air). The initial gas displacement is supersonic from the
outset: at $t = 1$~$\mu$s, the gas velocity is $v = a_\text{gas} \cdot
t = 8.16 \times 10^3$~m/s, which gives Mach number:
\begin{equation}
M = \frac{v}{c_s} = \frac{8160}{343} \approx 24.
\end{equation}

The flow is immediately \emph{highly} supersonic. A strong shock forms
at the leading edge of the displaced gas, compressing and heating it.
Behind the shock, the evacuated region has density limited by the
expansion wave speed. The net effect is that the near-field region
evacuates \emph{faster} than the simple ballistic estimate because
the supersonic expansion can create a near-vacuum behind the shock front.

For an ideal gas expanding into vacuum (Riemann problem), the density
behind an expansion wave drops exponentially:
\begin{equation}\label{eq:rho-expansion}
\rho(x, t) = \rho_0 \exp\!\left(-\frac{2}{(\gamma+1)}\frac{v(x,t)}{c_s}\right)
\end{equation}
where $\gamma = 1.4$ for air. Since $v/c_s \gg 1$ within a few
microseconds, the density drops to effectively zero in the near-field.

A more realistic profile including both the magnetic pressure drive and
the shock dynamics gives:
\begin{equation}\label{eq:rho-refined}
\rho_\text{near}(t) \approx \rho_0 \exp\!\left(-\frac{t}{t_\text{clear}}\right)^2,
\qquad t_\text{clear} \approx 5 \; \mu\text{s}.
\end{equation}

%=============================================================================
\section{Threshold Crossing Dynamics}\label{sec:threshold-crossing}
%=============================================================================

\subsection{When Does $\eta$ Cross $\etac$?}

The time-dependent threshold ratio in the evacuating region:
\begin{equation}
\eta(t) = \frac{U_\text{EM}}{\rho(t) \, c^2}
= \frac{B^2/(2\mu_0)}{\rho(t) \, c^2}.
\end{equation}

For $B = 50$~T, $U_\text{EM} = 9.95 \times 10^8$~J/m$^3$. The
threshold condition $\eta \geq \etac$ requires:
\begin{equation}
\rho(t) \leq \frac{U_\text{EM}}{\etac \, c^2}
= \frac{9.95 \times 10^8}{1.82 \times 10^{-3} \times 9 \times 10^{16}}
= \frac{9.95 \times 10^8}{1.64 \times 10^{14}}
= 6.07 \times 10^{-6} \; \si{kg/m^3}.
\end{equation}

The density must drop by a factor of:
\begin{equation}
\frac{\rho_\text{air}}{\rho_\text{thresh}} = \frac{1.225}{6.07 \times 10^{-6}}
= 2.02 \times 10^{5}.
\end{equation}

Using the Gaussian evacuation model (Eq.~\ref{eq:rho-refined}):
\begin{equation}
\rho_0 \exp\!\left(-\frac{t^2}{t_\text{clear}^2}\right) = \frac{\rho_0}{2.02 \times 10^5}
\end{equation}
\begin{equation}
\frac{t_\text{cross}^2}{t_\text{clear}^2} = \ln(2.02 \times 10^5) = 12.22
\end{equation}
\begin{equation}\label{eq:t-cross}
\boxed{
t_\text{cross} = t_\text{clear}\sqrt{12.22}
= 5 \times 3.50 = 17.5 \; \mu\text{s}.
}
\end{equation}

\paragraph{Result:} The threshold is first crossed approximately
\textbf{17.5~$\mu$s} after pulse initiation.

\subsection{Density at Threshold Crossing}

At $t = t_\text{cross}$:
\begin{equation}
\rho(t_\text{cross}) = 6.07 \times 10^{-6} \; \si{kg/m^3}
\approx 5 \times 10^{-6} \rho_\text{air}.
\end{equation}

This corresponds to a pressure of approximately 0.5~Pa---comparable to
a moderate vacuum ($\sim 4 \times 10^{-3}$~torr). The magnetic pressure
has transiently created near-vacuum conditions at sea level.

\subsection{Mass in the Above-Threshold Region at First Crossing}

At $t_\text{cross}$, the above-threshold region is the near-field zone
where the density has dropped below $\rho_\text{thresh}$. For a
magnet bore of radius $r_\text{bore} = 0.1$~m and an above-threshold
shell of thickness $\delta r \sim L = 0.1$~m:
\begin{equation}
V_\text{above} = 4\pi r_\text{bore}^2 \delta r
= 4\pi (0.1)^2 (0.1) = 1.26 \times 10^{-2} \; \si{m^3}.
\end{equation}

The mass of gas in this region at threshold crossing:
\begin{equation}\label{eq:mass-residual}
m_\text{gas} = \rho_\text{thresh} \times V_\text{above}
= 6.07 \times 10^{-6} \times 1.26 \times 10^{-2}
= 7.6 \times 10^{-8} \; \si{kg}
\approx 76 \; \si{\mu g}.
\end{equation}

\subsection{Gas Velocity Under $\kappa$-Coupling}

Once $\eta > \etac$, the $\kappa$-channel activates. The
DFD acceleration in the above-threshold region (from Agent~4,
Eq.~(\ref*{eq:a-MOND-like})) is:
\begin{equation}
a_\text{DFD}(r) = \sqrt{4\pi G \, U_\text{EM}(r)}.
\end{equation}

However, this is the MOND-like acceleration from the \emph{gravitational}
sourcing of $\psi$. The direct EM-$\psi$ coupling above threshold
generates a much stronger effect. The relevant acceleration is from the
full $\kappa$-enhanced $\psi$ gradient:
\begin{equation}
a_\kappa = \frac{c^2}{2}\frac{d\psi_\kappa}{dr}
\end{equation}
where $\psi_\kappa = \ka(\eta - \etac)$ in the above-threshold region.

For the field gradient over length $L$:
\begin{equation}
\frac{d\psi_\kappa}{dr} \sim \frac{\ka \, \eta}{L}
= \frac{51.4}{L} \cdot \frac{U_\text{EM}}{\rho c^2}.
\end{equation}

At threshold crossing ($\eta \approx \etac$):
\begin{equation}
a_\kappa \approx \frac{c^2}{2} \cdot \frac{\ka \, \etac}{L}
= \frac{(3 \times 10^8)^2}{2} \cdot \frac{51.4 \times 1.82 \times 10^{-3}}{0.1}
= 4.5 \times 10^{16} \cdot 0.936
= 4.2 \times 10^{16} \; \si{m/s^2}.
\end{equation}

This is an extraordinarily large acceleration. However, it applies
only to the residual gas ($76$~$\mu$g) and only for the brief
above-threshold window. As the gas is expelled, $\eta$ increases
further (positive feedback), but the gas density drops, so the
accelerated mass decreases.

The above-threshold window duration (from threshold crossing at
$t_\text{cross} = 17.5$~$\mu$s to pulse end at $t_\text{pulse} \sim
100$~$\mu$s):
\begin{equation}
\Delta t_\text{active} = t_\text{pulse} - t_\text{cross}
\approx 100 - 17.5 = 82.5 \; \mu\text{s}.
\end{equation}

The velocity acquired by the residual gas:
\begin{equation}\label{eq:v-gas}
v_\text{gas} = a_\kappa \cdot \Delta t_\text{active}
= 4.2 \times 10^{16} \times 82.5 \times 10^{-6}
= 3.5 \times 10^{12} \; \si{m/s}.
\end{equation}

This exceeds $c$, which means the gas immediately reaches relativistic
velocities and we must cap at a fraction of $c$. The DFD framework
operates within general relativity, so the actual terminal velocity
is limited by the energy budget. The kinetic energy available is:
\begin{equation}
E_\text{kin} = \ka \, U_\text{EM} \, V_\text{above}
= 51.4 \times 9.95 \times 10^8 \times 1.26 \times 10^{-2}
= 6.45 \times 10^8 \; \si{J}.
\end{equation}

For $m_\text{gas} = 7.6 \times 10^{-8}$~kg, the Lorentz factor:
\begin{equation}
(\gamma - 1)m_\text{gas}c^2 = E_\text{kin}
\end{equation}
\begin{equation}
\gamma - 1 = \frac{E_\text{kin}}{m_\text{gas}c^2}
= \frac{6.45 \times 10^8}{7.6 \times 10^{-8} \times 9 \times 10^{16}}
= \frac{6.45 \times 10^8}{6.84 \times 10^9} = 0.094.
\end{equation}

So $\gamma \approx 1.094$, giving:
\begin{equation}\label{eq:v-final}
\boxed{
v_\text{gas} = c\sqrt{1 - \gamma^{-2}}
= c\sqrt{1 - 0.836}
= c\sqrt{0.164}
\approx 0.41\,c
\approx 1.2 \times 10^8 \; \si{m/s}.
}
\end{equation}

\paragraph{Result:} Residual gas in the above-threshold region is
accelerated to $\sim 0.4\,c$---mildly relativistic. This is far more
efficient than any chemical or electric propulsion exhaust velocity.

\subsection{Reality Check on Energy Coupling}

The above assumes all the EM energy in the above-threshold volume
couples to the gas kinetic energy. In practice, the coupling efficiency
$\xi$ will be less than unity. The $\kappa$-channel describes a
\emph{modification of the refractive index}, which generates a
$\psi$-gradient that accelerates mass. The efficiency depends on the
geometry of the gradient and the gas distribution.

A conservative estimate uses $\xi = 0.01$ (1\% coupling):
\begin{equation}
(\gamma - 1) = 0.094 \times \xi = 9.4 \times 10^{-4},
\end{equation}
giving $v_\text{gas} \approx 0.043\,c \approx 1.3 \times 10^7$~m/s.

Even at 1\% efficiency, the exhaust velocity is $10^7$~m/s---three
orders of magnitude above the best ion thrusters.

%=============================================================================
\section{Impulse Per Pulse}\label{sec:impulse}
%=============================================================================

\subsection{Momentum Delivered}

The impulse per pulse is the momentum of the ejected gas:
\begin{equation}
I_\text{pulse} = \gamma \, m_\text{gas} \, v_\text{gas}.
\end{equation}

At full coupling ($\xi = 1$):
\begin{equation}
I_\text{pulse}^{(\xi=1)} = 1.094 \times 7.6 \times 10^{-8} \times 1.2 \times 10^8
= 1.094 \times 9.12 \times 10^{0}
\approx 10.0 \; \si{N \cdot s}.
\end{equation}

At conservative coupling ($\xi = 0.01$):
\begin{equation}
I_\text{pulse}^{(\xi=0.01)} = 1.001 \times 7.6 \times 10^{-8} \times 1.3 \times 10^7
= 9.9 \times 10^{-1}
\approx 1.0 \; \si{N \cdot s}.
\end{equation}

\subsection{Integral Formulation}

More precisely, the impulse is:
\begin{equation}
I = \int_0^{t_\text{pulse}} F(t) \, dt
= \int_{t_\text{cross}}^{t_\text{pulse}} \ka \frac{U_\text{EM}(t)}{c^2} \cdot a_\kappa(t) \cdot V_\text{above}(t) \, dt.
\end{equation}

The above-threshold volume grows as the gas evacuates (the region where
$\eta > \etac$ expands as density drops). However, the gas mass in the
growing volume also decreases. The dominant contribution comes from the
first few microseconds after threshold crossing, when there is still
appreciable mass in the newly-above-threshold region.

%=============================================================================
\section{Average Thrust at Various Repetition Rates}\label{sec:thrust}
%=============================================================================

The average thrust is:
\begin{equation}
\langle F \rangle = I_\text{pulse} \times f_\text{rep}.
\end{equation}

\begin{center}
\begin{tabular}{cccc}
\toprule
$f_\text{rep}$ (Hz) & $\langle F \rangle$ ($\xi = 1$) & $\langle F \rangle$ ($\xi = 0.01$) & Notes \\
\midrule
1 & 10 N & 1.0 N & Slow pulsing \\
10 & 100 N & 10 N & Moderate \\
100 & 1 kN & 100 N & High rate \\
1{,}000 & 10 kN & 1 kN & Fast pulsing \\
10{,}000 & 100 kN & 10 kN & Very fast \\
\bottomrule
\end{tabular}
\end{center}

\paragraph{Assessment:} Even at the conservative 1\% coupling, 100~Hz
repetition gives 100~N of thrust---enough to lift $\sim$10~kg against
gravity. At 1~kHz, 1~kN lifts $\sim$100~kg. To lift a vehicle of mass
$M$, we need:
\begin{equation}
f_\text{rep} = \frac{Mg}{I_\text{pulse}}
= \frac{M \times 9.81}{I_\text{pulse}}.
\end{equation}

For a 10{,}000~kg vehicle with $\xi = 0.01$:
\begin{equation}
f_\text{rep} = \frac{10^4 \times 9.81}{1.0} \approx 10^5 \; \si{Hz} = 100 \; \si{kHz}.
\end{equation}

This is a challenging but not impossible repetition rate for pulsed
magnets.

%=============================================================================
\section{Power Requirements}\label{sec:power}
%=============================================================================

\subsection{Energy Per Pulse: Full 50~T Field}

The magnetic energy stored in a volume $V$ at field $B$:
\begin{equation}
E_\text{pulse} = \frac{B^2}{2\mu_0} \cdot V.
\end{equation}

For a 50~T field in a volume $V = 10^{-3}$~m$^3$ (1 litre):
\begin{equation}
E_\text{pulse} = \frac{(50)^2}{2 \times 4\pi \times 10^{-7}} \times 10^{-3}
= 9.95 \times 10^8 \times 10^{-3}
= 9.95 \times 10^5 \; \si{J}
\approx 1 \; \si{MJ}.
\end{equation}

\subsection{Average Power}

\begin{center}
\begin{tabular}{ccc}
\toprule
$f_\text{rep}$ (Hz) & $P_\text{avg}$ (W) & Feasibility \\
\midrule
1 & 1 MW & Large generator \\
10 & 10 MW & Nuclear reactor \\
100 & 100 MW & Large reactor \\
1{,}000 & 1 GW & Unrealistic \\
10{,}000 & 10 GW & Unachievable \\
\bottomrule
\end{tabular}
\end{center}

Full 50~T pulsing at the rates needed for vehicle-scale thrust
($\sim$kHz) requires GW-class power---not feasible with current
technology.

\subsection{Energy Efficiency}

The thrust-to-power ratio at $\xi = 0.01$:
\begin{equation}
\frac{\langle F \rangle}{P_\text{avg}}
= \frac{I_\text{pulse} \times f_\text{rep}}{E_\text{pulse} \times f_\text{rep}}
= \frac{I_\text{pulse}}{E_\text{pulse}}
= \frac{1.0}{10^6}
= 10^{-6} \; \si{N/W}.
\end{equation}

This is poor compared to electric propulsion ($\sim 0.01$~N/W for
ion thrusters) but the exhaust velocity is five orders of magnitude
higher. The specific impulse:
\begin{equation}
I_\text{sp} = \frac{v_\text{exhaust}}{g_0}
= \frac{1.3 \times 10^7}{9.81}
\approx 1.3 \times 10^6 \; \si{s}.
\end{equation}

This is extraordinary---six orders of magnitude beyond chemical
rockets and three beyond the best ion thrusters.

%=============================================================================
\section{Optimised Pulsed Operation}\label{sec:optimised}
%=============================================================================

\subsection{Concept: Modulate Rather Than Pulse from Zero}

The energy cost can be dramatically reduced by maintaining a persistent
base field $B_\text{low}$ (requiring no continuous power if using
superconducting coils) and pulsing only the increment $\Delta B$:
\begin{equation}
B(t) = B_\text{low} + \Delta B(t), \qquad
\Delta B(t) = (B_\text{high} - B_\text{low}) \cdot p(t),
\end{equation}
where $p(t)$ is a pulsed waveform.

The pulse energy is:
\begin{equation}
\Delta E = \frac{B_\text{high}^2 - B_\text{low}^2}{2\mu_0} \cdot V_\text{coil}
= \frac{(B_\text{high} + B_\text{low})(B_\text{high} - B_\text{low})}{2\mu_0} \cdot V_\text{coil}.
\end{equation}

\subsection{Design Point: $B_\text{low} = 18$~T, $B_\text{high} = 22$~T}

These values are chosen because:
\begin{itemize}
\item 18~T is achievable with persistent-mode high-temperature
  superconductors (REBCO tape at 4~K routinely achieves 20+~T).
\item The 4~T modulation $\Delta B = 4$~T is within pulsed coil
  capability at high repetition rates.
\item The 22~T peak field is well within non-destructive pulsed
  magnet limits.
\end{itemize}

The energy per pulse for $V_\text{coil} = 0.01$~m$^3$:
\begin{align}
\Delta E &= \frac{(22)^2 - (18)^2}{2 \times 4\pi \times 10^{-7}} \times 0.01 \notag \\
&= \frac{484 - 324}{2.513 \times 10^{-6}} \times 0.01 \notag \\
&= \frac{160}{2.513 \times 10^{-6}} \times 0.01 \notag \\
&= 6.37 \times 10^{7} \times 0.01 \notag \\
&= 6.37 \times 10^{5} \; \si{J} \approx 640 \; \si{kJ}.
\end{align}

\subsection{Threshold Analysis for 22~T}

At $B_\text{high} = 22$~T:
\begin{equation}
U_\text{EM}(22\,\text{T}) = \frac{(22)^2}{2\mu_0}
= \frac{484}{2.513 \times 10^{-6}}
= 1.93 \times 10^{8} \; \si{J/m^3}.
\end{equation}

The threshold density:
\begin{equation}
\rho_\text{thresh}(22\,\text{T})
= \frac{U_\text{EM}}{\etac \, c^2}
= \frac{1.93 \times 10^8}{1.82 \times 10^{-3} \times 9 \times 10^{16}}
= \frac{1.93 \times 10^8}{1.64 \times 10^{14}}
= 1.18 \times 10^{-6} \; \si{kg/m^3}.
\end{equation}

The density reduction factor required:
\begin{equation}
\frac{\rho_\text{air}}{\rho_\text{thresh}} = \frac{1.225}{1.18 \times 10^{-6}}
= 1.04 \times 10^{6}.
\end{equation}

This is a factor of $\sim$5 more demanding than the 50~T case. Using
the Gaussian evacuation model, the threshold crossing time scales as:
\begin{equation}
t_\text{cross}(22\,\text{T}) = t_\text{clear}(22\,\text{T}) \sqrt{\ln(1.04 \times 10^6)}
= t_\text{clear} \sqrt{13.85}.
\end{equation}

The magnetic pressure at 22~T:
\begin{equation}
P_\text{mag}(22\,\text{T}) = 1.93 \times 10^8 \; \si{Pa}
\approx 2000 \; \text{atm}.
\end{equation}

The initial acceleration:
\begin{equation}
a_\text{gas}(22\,\text{T}) = \frac{1.93 \times 10^8}{1.225 \times 0.1}
= 1.58 \times 10^{9} \; \si{m/s^2}.
\end{equation}

The clearing time:
\begin{equation}
t_\text{clear}(22\,\text{T}) = \sqrt{\frac{0.2}{1.58 \times 10^9}}
= \sqrt{1.27 \times 10^{-10}}
= 1.13 \times 10^{-5} \; \si{s}
\approx 11.3 \; \mu\text{s}.
\end{equation}

Threshold crossing:
\begin{equation}
t_\text{cross}(22\,\text{T}) = 11.3 \times \sqrt{13.85}
= 11.3 \times 3.72 = 42 \; \mu\text{s}.
\end{equation}

\subsection{Impulse at 22~T}

The available EM energy in the above-threshold volume at 22~T is lower
than at 50~T, and threshold crossing takes longer. The residual gas
mass at crossing:
\begin{equation}
m_\text{gas}(22\,\text{T}) = 1.18 \times 10^{-6} \times 1.26 \times 10^{-2}
= 1.49 \times 10^{-8} \; \si{kg}.
\end{equation}

The available coupling energy ($\xi = 0.01$):
\begin{equation}
E_\text{couple} = \xi \, \ka \, U_\text{EM} \, V_\text{above}
= 0.01 \times 51.4 \times 1.93 \times 10^8 \times 1.26 \times 10^{-2}
= 1.25 \times 10^6 \; \si{J}.
\end{equation}

Exhaust velocity (non-relativistic regime at $\xi = 0.01$):
\begin{equation}
v_\text{exhaust} = \sqrt{\frac{2 E_\text{couple}}{m_\text{gas}}}
= \sqrt{\frac{2 \times 1.25 \times 10^6}{1.49 \times 10^{-8}}}
= \sqrt{1.68 \times 10^{14}}
= 1.30 \times 10^7 \; \si{m/s}
\approx 0.043\,c.
\end{equation}

The impulse per pulse:
\begin{equation}
I_\text{pulse}(22\,\text{T}) = m_\text{gas} \times v_\text{exhaust}
= 1.49 \times 10^{-8} \times 1.30 \times 10^7
= 0.19 \; \si{N \cdot s}.
\end{equation}

\subsection{Thrust and Power at 100~Hz}

\begin{equation}
\langle F \rangle = I_\text{pulse} \times f_\text{rep}
= 0.19 \times 100 = 19 \; \si{N}.
\end{equation}

\begin{equation}
P_\text{avg} = \Delta E \times f_\text{rep}
= 6.37 \times 10^5 \times 100
= 6.37 \times 10^7 \; \si{W}
\approx 64 \; \si{MW}.
\end{equation}

Thrust-to-power ratio:
\begin{equation}
\frac{\langle F \rangle}{P} = \frac{19}{6.37 \times 10^7}
= 3.0 \times 10^{-7} \; \si{N/W}.
\end{equation}

\subsection{Scaling Summary for Optimised Operation}

\begin{center}
\begin{tabular}{lcc}
\toprule
Parameter & 50~T full pulse & 18/22~T modulated \\
\midrule
$E_\text{pulse}$ & 1 MJ & 640 kJ \\
$t_\text{cross}$ & 17.5~$\mu$s & 42~$\mu$s \\
$m_\text{gas}$ per pulse & 76~$\mu$g & 15~$\mu$g \\
$v_\text{exhaust}$ ($\xi = 0.01$) & $0.043\,c$ & $0.043\,c$ \\
$I_\text{pulse}$ ($\xi = 0.01$) & 1.0 N$\cdot$s & 0.19 N$\cdot$s \\
$\langle F \rangle$ at 100 Hz & 100 N & 19 N \\
$P_\text{avg}$ at 100 Hz & 100 MW & 64 MW \\
$I_\text{sp}$ & $1.3 \times 10^6$ s & $1.3 \times 10^6$ s \\
\bottomrule
\end{tabular}
\end{center}

%=============================================================================
\section{Existing Pulsed Magnet Technology}\label{sec:technology}
%=============================================================================

\subsection{Current State of the Art}

\begin{center}
\begin{tabular}{llll}
\toprule
Facility/Technology & $B_\text{max}$ (T) & Duration & Notes \\
\midrule
Los Alamos NHMFL (destructive) & 100 & 10 ms & Single shot \\
Los Alamos NHMFL (non-dest.) & 65 & seconds & Repeated \\
Capacitor-driven pulsed & 50 & $\sim$100~$\mu$s & Routine \\
Magnetic forming (repetitive) & 30--40 & $\sim$10~$\mu$s & kHz capable \\
REBCO persistent (4~K) & 20+ & steady state & Zero power \\
Tokamak coils (ITER) & 13 & steady state & SC technology \\
\bottomrule
\end{tabular}
\end{center}

\subsection{Feasibility Assessment}

The key requirements for pulsed $\psi$-bubble operation:

\begin{enumerate}
\item \textbf{Base field (18~T persistent):} Achievable with
  state-of-the-art REBCO high-temperature superconducting tape operating
  at $\sim$4--20~K. The 32~T all-superconducting magnet at the NHMFL
  demonstrates that 18~T persistent fields are well within reach.
  Power consumption: effectively zero (persistent mode).

\item \textbf{Pulse increment (4~T modulation):} This requires pulsing
  a copper or REBCO insert coil from 0 to 4~T in $\sim$10~$\mu$s.
  The inductance of a small coil ($V = 0.01$~m$^3$, $L \sim 0.1$~m)
  is approximately:
  \begin{equation}
  \mathcal{L} \sim \mu_0 n^2 V \sim 10^{-4} \; \si{H}
  \end{equation}
  for $n \sim 10^4$~turns/m. The required current change:
  \begin{equation}
  \Delta I = \frac{\Delta B}{\mu_0 n} = \frac{4}{4\pi \times 10^{-7} \times 10^4}
  \approx 300 \; \si{A}.
  \end{equation}
  The voltage for 10~$\mu$s rise time:
  \begin{equation}
  V = \mathcal{L} \frac{\Delta I}{\Delta t}
  = 10^{-4} \times \frac{300}{10^{-5}} = 3000 \; \si{V}.
  \end{equation}
  This is entirely within the capability of modern capacitor banks
  and IGBT switching systems.

\item \textbf{Repetition rate (100~Hz):} For magnetic forming
  applications, repetition rates of 10--100~Hz at 30+~T are standard
  industrial practice. The coil heating is the main limitation, requiring
  active cooling between pulses.

\item \textbf{Power supply (64~MW):} This is the most challenging
  requirement. Options:
  \begin{itemize}
  \item \textbf{Compact nuclear reactor:} Modern small modular
    reactors (SMRs) deliver 60--300~MW$_\text{e}$. A single SMR
    could power the system.
  \item \textbf{Capacitor bank with reactor charging:} A 640~kJ
    capacitor bank charged at 100~Hz requires 64~MW average input.
    High-energy capacitor banks of this scale exist for fusion
    research (e.g., the NIF laser system stores 422~MJ).
  \item \textbf{Flywheel energy storage:} Could buffer between a
    lower-power source and the pulsed demand.
  \end{itemize}

\item \textbf{Thermal management:} At 64~MW with (optimistically)
  50\% electrical efficiency, 32~MW of waste heat must be rejected.
  For an atmospheric craft, radiative and convective cooling is
  possible.
\end{enumerate}

%=============================================================================
\section{Enhanced Designs}\label{sec:enhanced}
%=============================================================================

\subsection{Multiple Pulse Units}

An array of $N$ smaller pulse units, each operating at lower energy
but staggered in time, provides:
\begin{itemize}
\item Quasi-continuous thrust (ripple decreases as $1/N$).
\item Reduced per-unit energy ($E/N$ per coil).
\item Redundancy and failure tolerance.
\item Directional control by varying pulse amplitude across the array.
\end{itemize}

For $N = 100$ units at 10~Hz each (total effective rate 1~kHz):
\begin{equation}
\langle F \rangle_\text{total} = 100 \times 0.19 \times 10 = 190 \; \si{N},
\end{equation}
with each unit handling only 6.4~kJ/pulse at 10~Hz = 64~kW per unit.

\subsection{Shaped Field Geometry}

Instead of a spherical field, a \textbf{toroidal} or \textbf{conical}
field geometry can direct the gas ejection preferentially in one
direction, increasing the thrust-to-power ratio. A conical geometry
with half-angle $\theta_\text{cone} = 30^\circ$ concentrates the
impulse along the axis:
\begin{equation}
F_\text{directed} = \frac{F_\text{isotropic}}{1 - \cos\theta_\text{cone}}
= \frac{F_\text{isotropic}}{0.134}
\approx 7.5 \, F_\text{isotropic}.
\end{equation}

\subsection{Pre-Evacuation with Mechanical Pumping}

If the operating region is partially enclosed and pumped to moderate
vacuum ($\sim$100~Pa, readily achievable with roughing pumps), the
threshold field requirement drops by a factor of $\sim$100 in
density:
\begin{equation}
B_\text{thresh}(100\,\text{Pa}) = B_\text{thresh}(\text{atm})
\times \left(\frac{100}{10^5}\right)^{1/2} \approx 0.032 \times B_\text{thresh}.
\end{equation}

At 100~Pa ambient, a 22~T field would need only:
\begin{equation}
\frac{\rho_\text{thresh}}{\rho_{100\,\text{Pa}}} = \frac{1.18 \times 10^{-6}}{1.2 \times 10^{-3}}
= 10^{-3},
\end{equation}
requiring only a factor of 1{,}000 density reduction---achievable in
$\sim$5~$\mu$s even with modest magnetic pressure.

%=============================================================================
\section{Complete Parameter Summary}\label{sec:summary}
%=============================================================================

\begin{center}
\begin{tabular}{lll}
\toprule
\textbf{Parameter} & \textbf{Value} & \textbf{Source} \\
\midrule
\multicolumn{3}{c}{\textit{DFD Constants}} \\
$\ka = 3/(8\alpha)$ & 51.4 & Gauge emergence \\
$\etac = \alpha/4$ & $1.82 \times 10^{-3}$ & Threshold condition \\
\midrule
\multicolumn{3}{c}{\textit{Atmospheric Conditions}} \\
$\rho_\text{air}$ (sea level) & 1.225 kg/m$^3$ & Standard atmosphere \\
$U_\text{thresh}$ (sea level) & $2.01 \times 10^{14}$ J/m$^3$ & $\etac \rho c^2$ \\
$B_\text{thresh}$ (sea level) & 22{,}500 T & Static threshold \\
\midrule
\multicolumn{3}{c}{\textit{50~T Full Pulse}} \\
$P_\text{mag}$ & $10^9$ Pa ($10^4$ atm) & $B^2/(2\mu_0)$ \\
$a_\text{gas}$ & $8.2 \times 10^9$ m/s$^2$ & $P/(\rho L)$ \\
$t_\text{clear}$ (0.1 m) & 5.0 $\mu$s & $\sqrt{2L/a}$ \\
$t_\text{cross}$ & 17.5 $\mu$s & Gaussian model \\
$\rho$ at crossing & $6 \times 10^{-6}$ kg/m$^3$ & Threshold condition \\
$m_\text{gas}$ per pulse & 76 $\mu$g & $\rho V$ \\
$v_\text{exhaust}$ ($\xi = 0.01$) & $0.043\,c$ & Energy balance \\
$I_\text{pulse}$ ($\xi = 0.01$) & 1.0 N$\cdot$s & $mv$ \\
$E_\text{pulse}$ & 1 MJ & $B^2V/(2\mu_0)$ \\
\midrule
\multicolumn{3}{c}{\textit{18/22~T Optimised}} \\
$\Delta E$ per pulse & 640 kJ & $(B_h^2 - B_l^2)V/(2\mu_0)$ \\
$t_\text{cross}$ & 42 $\mu$s & Gaussian model \\
$m_\text{gas}$ per pulse & 15 $\mu$g & Lower density \\
$I_\text{pulse}$ ($\xi = 0.01$) & 0.19 N$\cdot$s & $mv$ \\
$I_\text{sp}$ & $1.3 \times 10^6$ s & $v_\text{ex}/g$ \\
$\langle F \rangle$ at 100 Hz & 19 N & $I \times f$ \\
$P_\text{avg}$ at 100 Hz & 64 MW & $\Delta E \times f$ \\
\bottomrule
\end{tabular}
\end{center}

%=============================================================================
\section{Conclusions and Next Steps}\label{sec:conclusions}
%=============================================================================

\subsection{Key Findings}

\begin{enumerate}
\item \textbf{Pulsed operation is physically viable.} A 50~T pulse
  evacuates air from the near-field region in $\sim$5~$\mu$s, creating
  transient near-vacuum conditions sufficient for threshold crossing
  at $\sim$17.5~$\mu$s.

\item \textbf{Extreme exhaust velocity.} The $\kappa$-channel
  accelerates residual gas to $\sim 0.04\,c$ even at 1\% coupling
  efficiency, yielding $I_\text{sp} \sim 10^6$~s---six orders of
  magnitude beyond chemical rockets.

\item \textbf{Power is the bottleneck.} The energy per pulse
  ($\sim$0.6--1~MJ) dominates the engineering challenge. At 100~Hz
  repetition, average power is $\sim$64~MW for the optimised 18/22~T
  scheme.

\item \textbf{Nuclear-powered atmospheric craft.} A compact reactor
  (SMR class, 60--300~MW$_e$) could power the system, enabling a
  DFD-propelled atmospheric vehicle with:
  \begin{itemize}
  \item No propellant consumption (uses ambient air as working fluid).
  \item Specific impulse $10^6$~s.
  \item Unlimited operational endurance (reactor limited).
  \item No moving parts in the propulsion system.
  \end{itemize}

\item \textbf{The coupling efficiency $\xi$ is the critical unknown.}
  The difference between $\xi = 1$ and $\xi = 0.01$ is the difference
  between a 19~N thruster and a 10~kN thruster at the same power.
  Determining $\xi$ experimentally is the highest-priority next step.

\item \textbf{Technology readiness.} Every component---REBCO
  superconductors, pulsed copper coils, capacitor banks, power
  electronics---exists today. The novel element is the DFD physics
  itself, which must be experimentally validated.
\end{enumerate}

\subsection{Experimental Roadmap}

\begin{enumerate}
\item \textbf{Phase 1: Threshold validation.} Operate a 20+~T pulsed
  magnet in a low-pressure chamber ($\sim$1~Pa) and search for anomalous
  gas acceleration signatures (the $\kappa$-channel effect).

\item \textbf{Phase 2: Atmospheric pulse test.} Operate in ambient air
  and measure:
  \begin{itemize}
  \item Time-resolved density profile (interferometry).
  \item Threshold crossing signature (optical emission from accelerated gas).
  \item Impulse measurement (force transducer on magnet mount).
  \end{itemize}

\item \textbf{Phase 3: Repetitive operation.} Demonstrate sustained
  average thrust at 10--100~Hz repetition with integrated force
  measurement.

\item \textbf{Phase 4: Optimised geometry.} Implement conical field
  shaping and multi-unit arrays for directed thrust demonstration.
\end{enumerate}

\appendix

%=============================================================================
\section{Derivation: Gas Return Time}\label{app:return}
%=============================================================================

After the pulse ends, the evacuated region refills by rarefaction wave
propagation at the speed of sound ($c_s = 343$~m/s for air):
\begin{equation}
t_\text{refill} = \frac{L_\text{evacuated}}{c_s}
= \frac{0.1}{343} \approx 290 \; \mu\text{s}.
\end{equation}

At 100~Hz repetition (10~ms between pulses), there is ample time
($\sim$35 refill times) for complete density recovery between pulses.
At 10~kHz (100~$\mu$s between pulses), the density may not fully
recover, creating a partial vacuum that actually \emph{reduces} the
threshold requirement for the next pulse---a potentially beneficial
resonance effect.

%=============================================================================
\section{Derivation: Relativistic Energy-Momentum}\label{app:relativistic}
%=============================================================================

For the exhaust gas with Lorentz factor $\gamma$, the four-momentum is:
\begin{equation}
p^\mu = (\gamma m c, \gamma m \mathbf{v}).
\end{equation}

The impulse per pulse (spatial component):
\begin{equation}
I = \gamma m v = \gamma m c \sqrt{1 - \gamma^{-2}}
= mc\sqrt{\gamma^2 - 1}.
\end{equation}

For $\gamma = 1.094$ ($\xi = 1$):
\begin{equation}
I = mc\sqrt{1.197 - 1} = mc\sqrt{0.197}
= mc \times 0.444
= 7.6 \times 10^{-8} \times 3 \times 10^8 \times 0.444
= 10.1 \; \si{N \cdot s}.
\end{equation}

For $\gamma \approx 1 + 9.4 \times 10^{-4}$ ($\xi = 0.01$):
\begin{equation}
I \approx m \sqrt{2(\gamma - 1)} \cdot c
= 7.6 \times 10^{-8} \times \sqrt{1.88 \times 10^{-3}} \times 3 \times 10^8
= 7.6 \times 10^{-8} \times 0.0434 \times 3 \times 10^8
\approx 1.0 \; \si{N \cdot s}.
\end{equation}

%=============================================================================
\section{Comparison with Other Propulsion Concepts}\label{app:comparison}
%=============================================================================

\begin{center}
\begin{tabular}{lccl}
\toprule
System & $I_\text{sp}$ (s) & Thrust/Power (N/W) & Status \\
\midrule
Chemical rocket & $250$--$450$ & $\sim 1$ & Mature \\
Ion thruster & $3000$--$10{,}000$ & $\sim 0.01$ & Flight proven \\
Hall thruster & $1500$--$3000$ & $\sim 0.02$ & Flight proven \\
VASIMR & $3000$--$30{,}000$ & $\sim 0.005$ & In development \\
Nuclear thermal & $800$--$1000$ & $\sim 0.5$ & Tested (1970s) \\
\textbf{DFD pulsed ($\xi=0.01$)} & $\mathbf{1.3 \times 10^6}$ & $\mathbf{3 \times 10^{-7}}$ & Theoretical \\
\textbf{DFD pulsed ($\xi=1$)} & $\mathbf{1.3 \times 10^6}$ & $\mathbf{1.6 \times 10^{-5}}$ & Theoretical \\
\bottomrule
\end{tabular}
\end{center}

The DFD pulsed system occupies a unique region of parameter space:
extremely high $I_\text{sp}$ with moderate thrust-to-power ratio. It is
most directly comparable to advanced electric propulsion concepts but
with three additional orders of magnitude in exhaust velocity. The
critical advantage for atmospheric operation is that it uses ambient
air as the working fluid, requiring no onboard propellant.

\end{document}
